
\subsubsection{6.1 Interpretation of Findings}

The study demonstrates a separation between measurement feasibility and construct validity. Linguistic markers were reliably extractable (H-E1 validated), but they did not validly measure agency preservation (H-M refuted across all criteria).

The low internal consistency ($\alpha =0.42)$ and weak inter-item correlations (r=0.22-0.35) suggest the three marker types tap different constructs rather than converging on a unified agency dimension. The consistent directional pattern (chosen < rejected) with negligible magnitude indicates potential confounding: chosen responses may be slightly more direct or concise for reasons related to response quality rather than agency preservation.

\subsubsection{6.2 Statistical Significance and Effect Size}

With N=169,352 pairs, the study demonstrates that p<0.001 can coexist with practically negligible effects (d=-0.018). The massive sample size provides >0.99 power to detect even trivial differences, making p-values uninformative for practical significance. The failure of the effect to replicate in the test split (d=-0.007, p=0.536), despite adequate power (0.95 for d=0.15), indicates that the train-split significance may be a statistical artifact of sample size rather than a robust phenomenon.

\subsubsection{6.3 Cross-Domain Transfer}

The findings indicate that linguistic markers validated in human psychology \citep{juanchich2017does} do not automatically transfer to AI-generated text in RLHF contexts. Human-authored language involves intentional communication where speakers deliberately use modal verbs to convey autonomy. AI-generated responses are produced through language model sampling optimized for preference signals, which may involve different generative processes.

\subsubsection{6.4 Limitations}

Results are specific to the HH-RLHF dataset (English, helpfulness-focused conversations). Generalization to other RLHF datasets, languages, or alignment methods was not tested. The study examined three specific marker types; alternative linguistic features were not explored. The analysis was conducted at the response level without stratification by conversation type or context, which may mask context-dependent patterns.

\subsubsection{6.5 Implications}

The negative result does not invalidate the bidirectional alignment framework \citep{shen2024bidirectional} or the importance of agency preservation. Rather, it demonstrates that computational operationalization cannot assume proxy validity based on cross-domain analogy. Future efforts to develop agency metrics may require: (1) direct user studies correlating linguistic features with human-annotated agency perception, (2) validation of alternative linguistic or behavioral features through multi-criterion testing, or (3) methods that move beyond linguistic proxies to behavioral or interaction-based measures.
